\documentclass[12pt, a4paper]{article}
\usepackage[utf8]{inputenc}
\usepackage[spanish]{babel}
\usepackage{amsmath, amssymb}
\usepackage{geometry}
\usepackage{enumitem}
\geometry{top=2.5cm, bottom=2.5cm, left=2.5cm, right=2.5cm}

\title{Resolución de Ejercicio: Técnicas de Conteo}
\author{}
\date{}

\begin{document}

\maketitle

\section*{Enunciado}
18. Disponemos de cinco cartas, cada una de ellas tiene grabada una letra: A, B, C, C y C. ¿De cuántas formas diferentes se pueden colocar en la mesa las cinco cartas, una al lado de la otra formando una hilera? Ejemplo: pueden estar colocadas de la siguiente forma ACBCC.

\vspace{0.5cm}
\hrule
\vspace{0.5cm}

\section*{Solución Paso a Paso}

\subsection*{1. Análisis e identificación del modelo}
Para determinar las distintas formas de colocar las cartas en una hilera, analizamos las características del ordenamiento:
\begin{itemize}[label=\checkmark]
    \item \textbf{¿Importa el orden?} Sí. Al formar una hilera, la posición relativa de cada carta genera una secuencia distinta (por ejemplo, ABC CC es diferente de CAB CC)[cite: 1].
    \item \textbf{¿Intervienen todos los elementos?} Sí. Se solicita colocar en la mesa exactamente las cinco cartas de las que disponemos ($m=5$)[cite: 1].
    \item \textbf{¿Existen elementos repetidos?} Sí. Dentro de nuestro conjunto total, existen elementos de la misma clase que son indistinguibles entre sí; específicamente, la letra ``C'' se repite tres veces[cite: 1].
\end{itemize}

Dado que importa el orden, intervienen todos los elementos del conjunto, pero hay elementos idénticos repetidos en la muestra, el modelo combinatorio adecuado es el de \textbf{Permutaciones con repetición}[cite: 1].

\subsection*{2. Planteo matemático y cálculo}
La fórmula estadística para calcular los ordenamientos de un conjunto de $m$ elementos donde hay $n_1, n_2, \dots, n_k$ elementos de distintas clases es[cite: 1]:
\[ P'_{n_1, n_2, \dots, n_k} = \frac{(n_1 + n_2 + \dots + n_k)!}{n_1! \times n_2! \times \dots \times n_k!} = \frac{m!}{n_1! \times n_2! \times \dots \times n_k!} \]

Definimos nuestra población total $m = 5$ cartas y desglosamos las repeticiones por clase:
\begin{itemize}
    \item Letra A: $n_1 = 1$
    \item Letra B: $n_2 = 1$
    \item Letra C: $n_3 = 3$
\end{itemize}

Procedemos a calcular, dividiendo el total de permutaciones posibles entre los ordenamientos internos de las cartas repetidas para no contabilizarlas como arreglos distintos[cite: 1]:
\begin{align*}
P'_{1,1,3} &= \frac{5!}{1! \times 1! \times 3!} \\
P'_{1,1,3} &= \frac{5 \times 4 \times 3!}{1 \times 1 \times 3!} \\
P'_{1,1,3} &= 5 \times 4 \\
P'_{1,1,3} &= 20
\end{align*}

\subsection*{3. Conclusión}
Existen 20 ordenamientos únicos posibles para alinear las cinco cartas dadas las repeticiones de la letra C.

\vspace{0.5cm}

\textbf{Respuesta:} Las cartas se pueden colocar de \textbf{20 formas diferentes}.

\end{document}
